\documentclass[11pt]{article}
\usepackage[a4paper,margin=23mm]{geometry}
\usepackage[T1]{fontenc}
\usepackage{lmodern,amsmath,amssymb,amsthm,mathtools,microtype,booktabs}
\usepackage[colorlinks=true,linkcolor=blue,citecolor=blue,urlcolor=blue]{hyperref}
\newcommand{\E}{\mathbb E}
\newcommand{\one}{\mathbf 1}
\newcommand{\norm}[1]{\left\lVert#1\right\rVert}
\newtheorem{theorem}{Theorem}
\newtheorem{lemma}{Lemma}
\title{An open-chord transfer for same-disorder curvature\\
\large Local source geometry and the Slater selected representation}
\author{Alan Day \and Chengheng\thanks{AI research assistant (OpenAI).}}
\date{8 November 2026\quad Local research draft}
\begin{document}\maketitle
\begin{abstract}
The fixed-Slater composite source formula is reduced to an explicit
recursion on open endpoints and two kinds of links. A normalized
two-label collision can retain only one quadratic mark, a pair of
nested quadratic marks, or linear marks on the same Gaussian edge.
This gives a complete source classification, including arbitrary
selected spectators. The transfer sums all fixed mixed words in the
Hamiltonian, size fluctuation, curvature and selected operators without
storing closed diagram history. Its reachable module carries the
positive Gram form inherited from the microscopic limit. The
auxiliary configurations themselves are not an orthonormal basis.
\end{abstract}

\section{Definitions and the source formula used as input}
Use standard CAR, the overlapping Hermitian complex Gaussian $p$-body
Hamiltonian, one common disorder $J$, and a fixed half-filled Slater
vector $O$. The state is annealed,
$\varphi_N=\E\langle O|\,\cdot\,|O\rangle$.
Take $m=N/2$, $\lambda_N=p^2/N\to\lambda>0$, at fixed word length.
Let $C_Q=\E H^2$ and $C_m$ be its scalar on charge $m$. Define
\begin{align}
 DX&=QX-\sum_i c_i^\dagger Xc_i,&
 \mathcal L X&=\tfrac12\bigl(NX-\sum_i(c_i^\dagger Xc_i+c_iXc_i^\dagger)\bigr),\nonumber\\
 h_N&=H_m/\sqrt{C_m},&
 W_N&=(p-\mathcal L-\lambda_N)H|_m/\sqrt{C_m},\nonumber\\
 j_N&=C_m^{-1}(2p-D-\kappa_N)(H^2-C_Q)|_m,&
 \kappa_N&=2\lambda_N\frac{1-q_N}{1+q_N},\quad q_N=e^{-4\lambda_N}.
\end{align}
Loss maps act on full Fock space before sector restriction. The
denominator is the fixed target-sector constant, not $C_Q$ inside a
loss map. Use dimensionless energy and $\hbar=1$. In the limit put
\begin{equation}
 q=e^{-4\lambda},\quad\gamma=e^{-\lambda},\quad
 \kappa=2\lambda\frac{1-q}{1+q},\quad
 W=\sqrt\lambda\,u,\quad
 v=\varphi(j^2)=4\lambda+\frac{16\lambda^2q}{1+q}.
\end{equation}
$M^\dagger$ is the normalized occupied-to-empty block of the same
original disorder; $\varphi(MM^\dagger)=1$.

The preceding project stage expands a $j$ into two adjacent $H$ slots
with a whole-block $D$ source, excluding its internal Gaussian pair.
$W$ is one $H$ slot with an $\mathcal L$ source. An edge is unrestricted
($G$), or forward selected ($F$). Backward edges vanish in this limit.
For a compatible pairing $\pi$, write $b$ for mixed $H$/selected edges
and $\nu$ for crossings plus nestings whose inner edge is $F$. Its weight is
\begin{equation}\label{input}
 \gamma^b[z_{\rm all}]
 \exp\!\left(\lambda K_\pi(z)
       -\lambda\sum_Wz_a-\kappa\sum_jz_a\right),
 \qquad K_\pi(0)=-4\nu .
\end{equation}
All source variables are commuting and square-zero. We reorganize
this already established fixed-word limit; no new $N$-limit is taken.

\newpage
\section{The complete local source classification}
Write $\pi(a)$ for the partner of endpoint $a$. A quadratic mark at
$(a,a+1)$ meets two different edges $e,f$. Let
$\epsilon=1-2\one_{e\text{ crosses }f}$. Define $\alpha_a$ as follows:
\begin{equation}\label{alpha}
 \alpha_a=\begin{cases}
 2\epsilon\one_{e_*\text{ is }G},&\pi(a),\pi(a+1)>a+1,
     \ e_*\text{ has the smaller right endpoint},\\
 2\epsilon\one_{e_*\text{ is }G},&\pi(a),\pi(a+1)<a,
     \ e_*\text{ has the larger left endpoint},\\
 2\epsilon,&\text{the partners lie on opposite sides.}
 \end{cases}
\end{equation}
\begin{lemma}\label{local}
The nonconstant coefficients of $K_\pi$ are precisely these:
\begin{itemize}
\item $K_W=1$ on a $G$ edge and $0$ on an $F$ edge;
$K_{WW'}=1$ if both marks are endpoints of the same $G$ edge.
\item $K_j=\alpha_a$. Two quadratic marks at $(a,a+1)$ and $(c,c+1)$,
$a<c$, have $K_{jj'}=4$ exactly when
$\pi(a)=c+1$, $\pi(a+1)=c$; otherwise it is zero.
\item Every mixed $K_{Wj}$ and every coefficient of source order at
least three vanishes.
\end{itemize}
\end{lemma}
\begin{proof}
The normalized collision expansion sums memberships containing two
labels. A site operator at an endpoint with neither label is $I$.
At a marked block containing just one label it is odd; both loss maps
multiply it by $1/2$. Its factor $1-z/2$ cancels the corresponding
singleton denominator, even when other marked blocks are present.
The Jordan--Wigner sign for that membership involves only its two
incident edges. Spectator edges therefore disappear identically.

If the labels are $P,Q$ of one edge, a quadratic block containing
one endpoint has $D(n)=n$, canceled by its two singleton factors
$(1-z/2)^2=1-z$. For an $\mathcal L$ mark,
$\mathcal L(n)=n-I/2$. Summing the empty and occupied matrix elements
with the factor $2$ in $K$ leaves the coefficients $1,1,1$ for the
two possible linear marks and their product. An $F$ edge forbids
this double membership, so all these coefficients vanish there.

For labels on different edges, every linear mark cancels. A surviving
$D$ mark must contain both incidences. Each edge has two endpoints,
so at most two such marks survive. Evaluation of the four choices
$P_eP_f,P_eQ_f,Q_eP_f,Q_eQ_f$ in the two site states gives the table
below. Two surviving marks consume all four endpoints: the same
calculation gives $4$ for the nested connection and $0$ for the crossed
one. Both edges then have $H$ endpoints, hence are $G$.
\end{proof}
After deleting spectators, label endpoints $0,1,2,3$ and order edges
by their left endpoints. The complete one-mark table is
\begin{center}
\begin{tabular}{ccrrrr}\toprule
Mark start&Pairing&$GG$&$GF$&$FG$&$FF$\\\midrule
0&$(02)(13)$&$-2$&$-2$&0&0\\
0&$(03)(12)$&2&0&2&0\\
1&$(01)(23)$&2&2&2&2\\
1&$(02)(13)$&$-2$&$-2$&$-2$&$-2$\\
2&$(02)(13)$&$-2$&0&$-2$&0\\
2&$(03)(12)$&2&0&2&0\\\bottomrule
\end{tabular}
\end{center}

\newpage
\section{Open configurations and single endpoints}
The cut-diagram transfer is an established method
\cite[Sec.~2.3]{BINT}; state-adapted dressed operators and their GNS
interpretation have related precedents \cite[Secs.~3.2--4]{RX}.
Here the additional information comes from the common-disorder sources.

Scan from right to left. A configuration $(s,\mathcal C)$ consists of
an ordered list $s$ of open right endpoints, oldest first, and disjoint
two-endpoint links $\mathcal C$. The five endpoint colors are
\begin{center}
\begin{tabular}{ccl}\toprule
Color&Edge type&Right endpoint\\\midrule
$H$&$G$&ordinary Hamiltonian slot\\
$\mathsf h$&$F$&ordinary Hamiltonian slot\\
$U$&$G$&linear marked slot $W$\\
$\mathsf w$&$F$&linear marked slot $W$\\
$P$&$F$&$M^\dagger$\\\bottomrule
\end{tabular}
\end{center}
Specifying $G/F$ when opening a Hamiltonian slot is a sum over future
compatible closures, not an additional factor. Every completed edge
has exactly one such type. The vacuum is the empty configuration.

Opening $H$ gives the two branches $H,\mathsf h$; opening $W$ gives
$U,\mathsf w$; $M^\dagger$ only appends $P$, and $M$ cannot open.
With $n$ open endpoints and a zero-based position $i$, the geometric
factor for closing is
\begin{equation}\label{closeq}
 \omega_i=\begin{cases}q^{n-1-i},&G,\\q^{n-1},&F.\end{cases}
\end{equation}
There are $n-1-i$ younger open edges, each crossing the closing edge.
When the closing edge is $F$, each of the $i$ older open edges also
forms an inner-$F$ nesting. This proves \eqref{closeq}.
Multiply it by the following endpoint coefficient; blanks are zero:
\begin{center}
\begin{tabular}{c|rrrrr}\toprule
Closing insertion&$H$&$\mathsf h$&$U$&$\mathsf w$&$P$\\\midrule
$H$&1&&&&$\gamma$\\
$W$&&&$\lambda$&&$-\lambda\gamma$\\
$M$&&$\gamma$&&$-\lambda\gamma$&1\\\bottomrule
\end{tabular}
\end{center}
Indeed an unpaired $W$ on a $G$ edge has coefficient
$\lambda(K_W-1)=0$. A pair of $W$ endpoints contributes $\lambda$;
a selected $W$ endpoint contributes $-\lambda$. This already sums
all linear sources, including their centering terms.

A link of type $S$ joins the two still-open edges of a singleton-assigned
quadratic mark. At its first closure multiply by
\begin{equation}\label{Sclose}
 2\lambda\epsilon\one_G-\kappa,\qquad
 \epsilon=\begin{cases}+1,&\text{younger linked endpoint closes first},\\
 -1,&\text{older linked endpoint closes first}.\end{cases}
\end{equation}
Delete that link; the other endpoint is now ordinary. A link of type
$R$ is reserved for a quadratic source pair and cannot close at an
individual endpoint. Relabel surviving positions after every deletion.

\newpage
\section{The explicit quadratic insertion}
Expand the coefficient of all quadratic sources in \eqref{input}
into singleton blocks of weight $\lambda\alpha-\kappa$ and pair blocks
of weight $4\lambda$. This is the finite set-partition expansion of
the square-zero exponential; no higher source block exists.

For a singleton-assigned $j$, process its right slot and then its left
slot as $H$. The two slots cannot pair with one another. There are four
possibilities, with all ordinary closure factors from the previous page:
\begin{description}
\item[OO.] Open both endpoints (each has its $G/F$ branches), join
them by an $S$ link, and defer the singleton weight to its first closure.
\item[OC.] Open the right endpoint, close an endpoint that was already
open before this insertion, and multiply by $-2\lambda-\kappa$.
The two incident edges cross.
\item[CO.] Close an old endpoint, open the left endpoint, and multiply
by $2\lambda-\kappa$. The two incident edges are disjoint.
\item[CC.] Close two distinct old endpoints with original indices
$i,k$, in that order. Put $\epsilon=+1$ if $i>k$ and $-1$ otherwise.
The younger original endpoint has index $\max(i,k)$. Multiply by
\begin{equation}\label{CC}
 2\lambda\epsilon\one_{s_{\max(i,k)}\text{ is }G}-\kappa .
\end{equation}
Any old $S$ link encountered during either closure is resolved as well.
Old $R$ endpoints are unavailable to these singleton branches.
\end{description}
The remaining branches assign the current quadratic mark to a source
pair:
\begin{description}
\item[Pair opening.] Append two $H$ endpoints and join them by an $R$
link, with coefficient one.
\item[Pair closing.] Choose an existing $R$ link $(a,b)$, $a<b$.
Close $b$ and then $a$, delete the link, and multiply by
\begin{equation}\label{Rclose}
 4\lambda\,q^{(n-1-b)+(n-2-a)}.
\end{equation}
The two edges are nested and unrestricted. Their crossings with
other open edges are precisely the two exponents in \eqref{Rclose}.
\end{description}
The links are always disjoint: an $S$ link is created only on two new
endpoints and is erased when either is removed; an $R$ link is created
and removed as a whole. Consequently the record stores no information
about already closed edges. Configurations with the same ordered colors
and links are merged by adding their coefficients.

For a prescribed finite word, one may discard configurations that
cannot close in its unprocessed prefix. The implementation checks the
remaining number of slots, $M$ closures for $\mathsf h,\mathsf w$,
$W$ closures for $U$, available $H$ slots for $H$, and $j$ insertions
for $R$ links. These are necessary conditions only; using them cannot
discard a valid completed path. This is exact pruning, not a numerical
cutoff or a bound on worst-case computational complexity.

\newpage
\section{Equivalence and the physical Gram form}
Let $T_A$ denote the preceding recursion for an insertion $A$, and
$\varepsilon$ the coefficient of the empty configuration. For a word
$w=A_1\cdots A_r$, set $T_w=T_{A_1}\cdots T_{A_r}$.
\begin{theorem}
For every fixed word in $h,W,j,M,M^\dagger$,
\begin{equation}\label{result}
 \varphi(w)=\varepsilon T_w\emptyset .
\end{equation}
The reachable span $V=\operatorname{span}\{T_P\emptyset\}$ has a
well-defined positive semidefinite form
\begin{equation}\label{Gram}
 G(T_P\emptyset,T_Q\emptyset)=\varphi(P^\dagger Q).
\end{equation}
Its null quotient is the joint polynomial GNS module, and the induced
actions satisfy $T_A^*=T_{A^\dagger}$ on that core.
\end{theorem}
\begin{proof}
Decorate a compatible forward pairing with one term of its finite
source-partition expansion. Each $j$ is either a singleton or paired
with another $j$ by the nested two-edge connection of Lemma~\ref{local}.
Scan the decorated pairing. Ordinary closures supply
$\gamma^bq^\nu$ and the linear-source weights. A singleton quadratic
mark follows exactly one OO, OC, CO or CC branch; the formula for
$\alpha$ gives respectively \eqref{Sclose}, the two split weights,
or \eqref{CC}. A source pair follows a pair opening and \eqref{Rclose}.
Conversely each completed path specifies this decorated pairing
uniquely. Thus the finite sums coincide. Merging identical open records
preserves all future weights, since those records contain every
unresolved source and edge. This proves \eqref{result}.

If $T_P\emptyset=0$, then for every polynomial $R$,
$\varphi(RP)=\varepsilon T_RT_P\emptyset=0$. Hence changes of
polynomial representative do not affect \eqref{Gram} in its second
slot, or in its first slot by Hermiticity. Positivity follows from the
established finite-$N$ limit of $\E\norm{P O}^2$. The form is the usual
positive moment form, whose null space is a left ideal by
Cauchy--Schwarz. Its quotient and adjoint relations therefore give
exactly the polynomial GNS construction.
\end{proof}
This argument gives the metric on the \emph{reachable} module. It does
not assign a Euclidean metric to formal open configurations, nor assert
that every such configuration is individually reachable. In practice a
word-basis Gram matrix is computed by applying the same recurrence to
$P_a^\dagger P_b$.

The previous polynomial-density theorem bounds the covariant algebra
on this entire GNS space. In particular
\begin{equation}
 B_h=\frac2{\sqrt{1-q}},\quad
 \norm W\le\sqrt\lambda B_h,\quad
 \norm j\le\frac{2\sqrt v}{\sqrt{1-q^4}}.
\end{equation}
For $S_K(t)=\sum_{r=0}^K(ith)^r/r!$, the limiting-space estimate is
\begin{equation}
 \norm{(e^{ith}-S_K(t))\psi}
 \le\norm\psi\,e^{|t|B_h}
              \frac{(|t|B_h)^{K+1}}{(K+1)!}.
\end{equation}
It provides controlled limiting covariant time calculations between
fixed preparations; it does not by itself exchange finite-$N$,
large-time and nonlinear-pinching limits.

\newpage
\section{Exact checks, scope and the next calculation}
The fresh checker uses rational arithmetic and imports only pure
functions from preceding stages. Its first completed scientific output
records:
\begin{center}
\begin{tabular}{lr}\toprule
Comparison&Count\\\midrule
Full local CAR source exponent versus geometric rule&2774\\
Selected source moments versus transfer (497 words, two points)&994\\
Old $H/U/P$ transfer: all six-letter $H,W,M,M^\dagger$ words&4096\\
Positive exact $14\times14$ Gram LDL decompositions&2\\
Unpruned versus pruned vacuum coefficient checks&5\\\bottomrule
\end{tabular}
\end{center}
The two formal points are $(\gamma,\lambda)=(1/2,1/3),(2/3,1/4)$.
They test rational identities off the physical curve
$\gamma=e^{-\lambda}$. Physical positivity is inherited from the
microscopic theorem, not inferred from these finite tests.

Longer examples provide an actual use of the transfer:
\begin{center}
\begin{tabular}{lrr}\toprule
Word&Expanded endpoints&Peak retained states\\\midrule
$h^{24}$&24&7\\
$j^8$&16&21\\
$j h^2 j W^2 j h^2 j$&14&10\\
$M j h^4 j M^\dagger$&10&18\\
$M^2 j h^2 j(M^\dagger)^2$&10&22\\\bottomrule
\end{tabular}
\end{center}
All are exact at the first formal point. The first reference is the
classical Touchard--Riordan formula \cite{JK},
\begin{equation}
 \varphi(h^{2n})=\frac1{(1-q)^n}\sum_{k=0}^n(-1)^k
 \left[\binom{2n}{n-k}-\binom{2n}{n-k-1}\right]q^{k(k+1)/2}.
\end{equation}
The next two use the preceding tracial mixed-$q$ law; the selected
examples use the original full source engine. A preliminary checker
was interrupted after mistakenly choosing full pairing enumeration for
the $h^{24}$ reference. It had written no scientific output. Replacing
that reference by the displayed classical sum allowed the first complete
run to finish. The interruption is retained in the diagnostics.

The new result is an arbitrary-fixed-word computational representation
with a physical Gram quotient. It does not classify the full selected
operator algebra, prove compactness of $M$ in this enlarged space, or
provide the spatial and mass generators of de Sitter holography. The
old bridge for actual curvature $a_N=(1-\mathcal P_H)x_N$, $x=j/\sqrt v$,
continues to apply to covariant blocks between fixed polynomial
preparations. General interlacing of nonlinear pinched blocks remains
a separate limit problem.

The next useful physical calculation is the full restriction of a
single selected excitation to the three-field covariant algebra,
followed by higher bilocal descendants. It can now be tested with
explicit ordered mixed moments while retaining the same original
disorder, before choosing a spectral or bulk interpretation.

\begin{thebibliography}{9}\small
\bibitem{BINT} M. Berkooz, M. Isachenkov, V. Narovlansky and G. Torrents,
\emph{Towards a full solution of the large N double-scaled SYK model},
\href{https://arxiv.org/abs/1811.02584}{arXiv:1811.02584}, Sec.~2.3.
\bibitem{RX} H. Rajgadia and J. Xu,
\emph{Emergent States and Algebras from the Double-Scaling limit of Pure States in SYK},
\href{https://arxiv.org/abs/2604.14387}{arXiv:2604.14387v3}, Secs.~3.2--4.
\bibitem{JK} M. Josuat-Verg\`es and J. S. Kim,
\emph{Touchard--Riordan formulas, T-fractions, and Jacobi's triple product identity},
\href{https://arxiv.org/abs/1101.5608}{arXiv:1101.5608};
FPSAC 2011 extended abstract, pp.~563--574.
\end{thebibliography}
\end{document}
